% Apendice A (revisao pos-EE). Gerado por 02_audit/build12_appendix.py; tabulars verbatim.
\begin{table}[htbp]\centering
\caption{Authorized capacity measured by offered places and by classes times the statutory module}
\label{tab:R_capcheck}
\footnotesize
\fitbox{%
\begin{tabular}{lrrrrrr}
\toprule
& \multicolumn{3}{c}{Offered places} & \multicolumn{3}{c}{Classes $\times$ statutory module} \\
\cmidrule(lr){2-4}\cmidrule(lr){5-7}
Age band & 2013S1 & 2019S1 & Change \% & 2013S1 & 2019S1 & Change \% \\
\midrule
Ber\c{c}\'ario I & 19,733 & 32,620 & 65.3 & 19,859 & 33,012 & 66.2 \\
Ber\c{c}\'ario II & 47,405 & 86,402 & 82.3 & 48,564 & 89,748 & 84.8 \\
Minigrupo I & 73,147 & 104,040 & 42.2 & 75,900 & 109,332 & 44.0 \\
Minigrupo II & 74,276 & 115,831 & 55.9 & 88,700 & 136,875 & 54.3 \\
\addlinespace
Total & 214,561 & 338,893 & 57.9 & 233,023 & 368,967 & 58.3 \\
\bottomrule
\end{tabular}}
\tabnote{\emph{Source:} municipal class profile (offered places, classes and enrolment by centre, age band and semester), 2013--2019. \emph{Sample:} first semesters of 2013 and 2019. The table compares the growth of two measures of authorized capacity: offered places, and the number of classes multiplied by the statutory module of the age band (7, 9, 12 or 25). The second measure is built only from class counts and statutory modules, so it is independent of enrolment by construction. Across band-semester cells, the ratio of offered places to module-based capacity lies between 0.837 and 0.994. The comparison is descriptive.}
\end{table}
